\section{Related Work}
\label{sec:related}

\paragraph{Benchmarks and measurements of agent-side injection.}
The measurement landscape has moved rapidly from static test sets toward generative and interactive protocols. AutoDojo generates novel injection scenarios programmatically and shows that many published defenses lose most of their advantage under unseen attack templates~\cite{ma2026autodojo}. StepJack extends the threat model to multi-step indirect injection against computer-use agents, where a single planted payload is consumed across several turns~\cite{zhan2026stepjack}; LivePI benchmarks agents in live, continuously mutating environments rather than frozen traces~\cite{zhao2026livepi}; NetInjectBench targets tool-using agents in network operations, where tickets and logs carry payloads~\cite{shayoni2026netinjectbench}; and IPI-proxy red-teams web-browsing agents through an intercepting proxy under domain whitelists~\cite{chiapei2026ipiproxy}. A parallel line measures what attacks actually achieve: covert injection designs aim to keep the user entirely unaware of a successful compromise, stressing the need for outcome-level rather than refusal-level metrics~\cite{lee2026will}. These build on the foundational static benchmarks InjecAgent and AgentDojo, which established that frontier models execute injected tool-directed instructions at material rates~\cite{zhan2024injecagent,debenedetti2024agentdojo}, and on automated attack generation, which showed that injection payloads need not be handcrafted~\cite{liu2024automatic,wang2025webinject,shi2025prompt,wang2025agentvigil,yang2026when}. \method{} complements this line at the measurement level: rather than scoring a single execution, our paired replay protocol measures \emph{post-compromise learning}---whether a defense extracted from earlier failures changes the outcome of a re-executed attack.

\paragraph{Defenses against indirect prompt injection.}
Defensive work spans four mechanisms. Preference- and structure-level methods reshape the model itself: StruQ enforces structured queries with a secure decoder, and SecAlign trains the model to prefer the user instruction over injected content~\cite{chen2024struq,chen2024secalign}. System-level designs constrain information flow at deployment time, as in AgentDojo's task-relevance checks and provable variants such as MELON~\cite{debenedetti2024agentdojo,zhu2025melon}. Runtime monitors inspect traffic around the agent: ClawGuard couples runtime detection with taint-based response, IPIGuard enforces a tool-dependency-graph execution discipline, and AttriGuard attributes unsafe tool invocations to their untrusted sources via causal analysis~\cite{zhao2026clawguard,an2025ipiguard,he2026attriguard}. Standalone classifiers detect injected content before it reaches the model~\cite{shaheer2025detecting}, while hidden-state probes flag injections from internal activations---though candidate evaluation protocols caution that high AUC alone overstates deployment readiness~\cite{li2026when}. Two structural gaps persist across families. First, defenses are evaluated within single episodes; none accumulates forensic evidence across episodes. Second, deployed agents continue to fail silently in production, and those failures are logged but rarely transformed into protection. \method{} targets exactly this residue: it treats flagged, already-completed executions as training material for a memory-based defense, and its experience boundary is built so that the defense never observes ground-truth labels---a property none of the above systems provides.

\paragraph{Agent memory and experience reuse.}
Experience reuse is becoming a first-class concern in agent systems. Recent surveys chart the shift from simple retrieval-augmented memory toward persistent, self-evolving memory systems for long-horizon agents~\cite{huang2026survey,zhou2026ready}, and retrieval-augmented agents demonstrably learn from accumulated experience to generalize to unseen tasks~\cite{ferraz2026retrievalaugmented}. The underlying mechanism---interleaving reasoning, acting, and observed feedback---dates to ReAct~\cite{yao2022react}, but its application to safety is sparse: memory content is typically harvested from successful task completions, not from security failures, and there is no selection mechanism deciding \emph{which} past episodes deserve to become memory. \method{} contributes the missing selection layer: a graph anomaly detector performs the triage, a weakly supervised scorer performs the localization, and a judge performs the admission decision, so that only localized, label-free, verdict-filtered evidence enters the memory store.

\paragraph{LLM-as-judge reliability.}
Our evaluation depends on automated verdicts, so judge reliability is a first-order design constraint. Recent work documents substantial failure modes: content-invariant wrappers flip safety-judge verdicts without changing substance~\cite{zhou2026style}, mechanistic analyses localize scoring bias inside judge models~\cite{xu2026inside}, policy-invariance tests expose verdicts that depend on asserted policies rather than behavior~\cite{weng2026beyond}, and even pinned judges at temperature zero exhibit verdict variance that forms a noise floor for leaderboards~\cite{ayyagari2026pinned}. For agent trajectories specifically, outcome-only judges miss silent process failures~\cite{mohammadi2026trajectoryjudge}, agent-as-judge pipelines evaluate intermediate steps but inherit the same model-family biases~\cite{zhuge2024agentasajudge}, and dedicated safety benchmarks judge risk awareness with mixed agreement~\cite{yuan2024rjudge}. Our protocol responds with two separations: the attack-success judge comes from a different model family than the agent and the synthesizer, and every judge verdict is auditable against the benchmark's recorded compromised execution, which anchors automated labels to recorded ground truth.

\paragraph{Graph learning for anomaly detection and explanation.}
Casting agent trajectories as graphs connects our detector to the graph-anomaly-detection literature, surveyed for time-series-linked settings by Ho et al.~\cite{ho2025graph}, with contrastive and multi-view formulations addressing label scarcity~\cite{liu2023revisiting}. Our detector is an edge-aware graph attention network~\cite{velikovi2017graph} in the relational-GCN lineage~\cite{schlichtkrull2017modeling} operating over frozen sentence embeddings~\cite{reimers2019sentencebert}, and its localization stage is a max-MIL node scorer in the spirit of GNNExplainer and PGExplainer, which produce subgraph explanations for graph-level predictions~\cite{ying2019gnnexplainer,luo2020parameterized}. The difference is purposeful: classical explainers optimize faithfulness to the model's decision, whereas our scorer is optimized to recover the compromised span \emph{as evidence for a downstream LLM judge}, so its objective trades explanatory faithfulness for evidence completeness under weak, graph-level-only supervision.
